% CA-TOSG v2 manuscript. V2-R42.
% Every number comes from tables/generated_numbers.tex, written by
% tools/build_v2_paper_numbers.py from the closed-out products. Do not type a figure by hand.
% The v1 manuscript (../main.tex) is FROZEN and is not touched by this file.
\documentclass[journal]{IEEEtran}
\usepackage{amsmath,amssymb}
\usepackage{graphicx}
\usepackage{booktabs}
\usepackage{multirow}
\usepackage{cite}
\usepackage{xcolor}
\usepackage{url}
% GENERATED by tools/build_v2_paper_numbers.py -- do not edit by hand (V2-R42 B-2).
% Every figure below traces to a closed-out product; see the publication manifest.
\newcommand{\TestRFFone}{0.87679}
\newcommand{\TestTauFone}{0.88178}
\newcommand{\TestRFPay}{0.00081}
\newcommand{\TestTauPay}{0.54102}
\newcommand{\TestDeltaFone}{-0.00499}
\newcommand{\TestLCB}{-0.00738}
\newcommand{\TestUCBPay}{-0.49987}
\newcommand{\TestSaving}{99.85}
\newcommand{\TestScenes}{16}
\newcommand{\TestRows}{47,740}
\newcommand{\TestNoCollab}{5.48}
\newcommand{\TestRhoE}{0.403}
\newcommand{\TestRhoL}{0.597}
\newcommand{\TestRhoF}{0.000}
\newcommand{\TestShortfall}{0.0024}
\newcommand{\TestRatio}{669}
\newcommand{\TestRhoEPct}{40.3}
\newcommand{\TestRhoLPct}{59.7}
\newcommand{\TestRhoFPct}{0.0}
\newcommand{\CulverRFFone}{0.87352}
\newcommand{\CulverTauFone}{0.87813}
\newcommand{\CulverRFPay}{0.00147}
\newcommand{\CulverTauPay}{0.49850}
\newcommand{\CulverDeltaFone}{-0.00461}
\newcommand{\CulverLCB}{-0.00541}
\newcommand{\CulverUCBPay}{-0.42272}
\newcommand{\CulverSaving}{99.70}
\newcommand{\CulverScenes}{4}
\newcommand{\CulverRows}{12,100}
\newcommand{\CulverNoCollab}{13.09}
\newcommand{\CulverRhoE}{0.407}
\newcommand{\CulverRhoL}{0.593}
\newcommand{\CulverRhoF}{0.000}
\newcommand{\CulverShortfall}{0.0004}
\newcommand{\CulverRatio}{339}
\newcommand{\CulverRhoEPct}{40.7}
\newcommand{\CulverRhoLPct}{59.3}
\newcommand{\CulverRhoFPct}{0.0}
\newcommand{\DeltaMargin}{0.005}
\newcommand{\TauStar}{16.5}
\newcommand{\CandIdx}{67}
\newcommand{\LambdaStar}{0.2}
\newcommand{\WtwoCTestAPfifty}{0.91708}
\newcommand{\WtwoCTestAPseventy}{0.82989}
\newcommand{\WtwoCTestRate}{5.2}
\newcommand{\WtwoCCulverAPfifty}{0.83195}
\newcommand{\WtwoCCulverAPseventy}{0.71588}
\newcommand{\WtwoCCulverRate}{9.1}


\begin{document}
\title{Task-Oriented Semantic Granularity Selection for\\Bandwidth-Constrained V2V Cooperative Perception}
\author{Peiyi~Yue}
\maketitle

\begin{abstract}
Cooperative perception must choose \emph{what} to transmit, not only how much to compress. We
formalise per-frame semantic granularity selection over three actions---ego-only (E), object-level
(L) and feature-level (F)---and evaluate it under a single detector, a single field of view and a
single ground truth, with communication charged through a measured LDPC/QAM chain rather than a
declared budget. The resulting method, CA-TOSG (channel-aware task-oriented semantic
granularity), reduced realised communication payload by \TestSaving\% on Test and
\CulverSaving\% on Culver-City relative to the frozen comparator. However, the preregistered
scene-level non-inferiority criterion was not met on either split, indicating that substantial
communication reduction was achieved without statistically establishing accuracy preservation.
Three mechanism results follow from the unified branch: partial recovery is necessary, since
all-or-nothing delivery of a $12{,}567$-codeword message is effectively dead at realistic loss
rates; \emph{where} a loss falls changes the task outcome even when \emph{how much} is lost is held
exactly fixed; and fusing a zeroed feature tensor is not equivalent to ego-only inference. In the
frozen policy the feature action is never selected, and we show this follows from the cost structure
and the sampled penalty grid rather than from a failure to learn.
\end{abstract}

\begin{IEEEkeywords}
Cooperative perception, V2V communication, semantic communication, bandwidth adaptation.
\end{IEEEkeywords}

\section{Introduction}
\IEEEPARstart{C}{ooperative} perception lets a vehicle see what its own sensors cannot. What it
should send is not settled: raw points, intermediate features and object lists differ by orders of
magnitude in size~\cite{hu2022where2comm,liu2020when2com} and in what they preserve. We treat the choice itself as the object of study.

Our contributions are three.
\begin{enumerate}
\item \textbf{Problem formalisation.} Per-frame granularity selection as a budget-constrained action
choice over $\{$E, L, F$\}$, with payload charged per frame through the transport chain the messages
actually traverse.
\item \textbf{A receiver-driven method.} The ego evaluates its own scene and channel and requests a
granularity; the selector's inputs are strictly ego-local and available \emph{before} the request.
\item \textbf{Experimental findings under a unified protocol.} One checkpoint, one field of view,
one ground truth, and a deterministic evaluation pipeline.
\end{enumerate}

The random forest used below is one instance of the selection layer, not the contribution. What is
being offered is the granularity-control formulation and the accounting protocol that makes a
payload claim checkable; any policy that consumes ego-local cues and emits a request slots into the
same framework, and we report below that this particular instance does not beat a tuned threshold.

\begin{figure*}[t]
\centering\includegraphics[width=\textwidth]{figures/fig1_system}
\caption{Data flow as implemented. The cue vector and the channel estimate are ego-side and
available before the request; the collaborator supplies message content only when asked. E issues no
message and therefore never enters the transport chain.}
\label{fig:system}
\end{figure*}

\section{Method}
\subsection{Action space and the receiver-driven request}
The ego selects one of three actions per frame and signals it with a two-bit request. That request is a fixed control
overhead, identical for all three actions, and is therefore not charged to any action's payload. E transmits nothing. L transmits the collaborator's detected boxes. F
transmits the collaborator's quantised bottleneck tensor.

\subsection{Communication accounting}
Payload is measured, not declared. For F the transmitted tensor is $739{,}200$ elements at
$w=8$~bits, packetised at $8{,}000$ payload bits with a $320$-bit header following common V2X framing
assumptions~\cite{ieee80211bd,3gpp37885}, giving $N_{cw}=12{,}567$
codewords and $B_F=3.14175$~Msym/frame. $B_F$ is derived from the codeword count rather than from
$(\text{info}+\text{header})/R/\log_2 M$, because the direct route ignores codeword padding.
L is charged per frame through the identical chain from its own box count.

\subsection{Transport}
The message is fragmented into $N_{cw}$ codewords; each succeeds or fails at the frame's SNR, and
the receiver reconstructs from what arrives, zero-filling failed blocks. All-or-nothing delivery is
retained as a pessimistic sensitivity arm. Three regimes are
evaluated: fragment-aware recovery, packet-level loss, and all-or-nothing message delivery
(Fig.~\ref{fig:recovery}).

\subsection{Selector inputs}
The cue vector has $23$ dimensions and is ego-local by construction: $19$ statistics of the ego's
own sweep, its own detected box count, a binary collaborator-availability flag, and two channel
fields. Every field is computable before the request is issued, so no input depends on a message the
ego has not yet asked for.

\section{Experimental Protocol}
All experiments use OPV2V~\cite{xu2022opv2v} with a PointPillars
detector~\cite{lang2019pointpillars}. One checkpoint serves all three actions; the field of view and ground truth are identical across
them. Evaluation is deterministic: the per-sample point shuffle is seeded on
$(\text{split},\text{scene},\text{frame},\text{cav})$, and two independent processes reproduce a run
bit for bit. The selector is fitted on the development split only and frozen---a random forest~\cite{breiman2001random}, candidate
\CandIdx{} at $\lambda=\LambdaStar$---before either held-out set is opened. The comparator is likewise
frozen: an SNR threshold at $\tau=\TauStar$, the best budget-feasible threshold at the primary cell.
The confirmatory unit is the scene, not the frame: frames within a
scene are consecutive samples of the same traffic and are far from independent. Intervals are
scene-level bootstraps over \TestScenes{} Test scenes and \CulverScenes{} Culver-City scenes,
covering \TestRows{} and \CulverRows{} evaluated rows respectively.

\section{Results}
\subsection{Primary comparison}
Table~\ref{tab:primary} gives the primary cell. On Test the payload falls by \TestSaving\% and on
Culver-City by \CulverSaving\%; on both splits the scene-level bootstrap lower bound on $\Delta F_1$
lies outside the preregistered margin of $-\DeltaMargin$, so accuracy non-inferiority is not established.
Both halves are stated together throughout: the saving is large---a factor of \TestRatio{} and
\CulverRatio{}---and the shortfall on the bound is \TestShortfall{} (Test) and \CulverShortfall{}
(Culver). The Culver-City interval rests on only \CulverScenes{} scenes and is correspondingly
coarse; it reproduces the shape of the Test result on much less scene diversity rather than
independently corroborating it.

\begin{table}[t]
\caption{Primary cell. $\Delta$ is CA-TOSG minus the frozen comparator; the criterion is defined on
the bootstrap bound, not the point estimate.}
\label{tab:primary}
\centering
\setlength{\tabcolsep}{3pt}\small
\begin{tabular}{lrrrr}
\toprule
 & \multicolumn{2}{c}{Test} & \multicolumn{2}{c}{Culver-City} \\
\cmidrule(lr){2-3}\cmidrule(lr){4-5}
 & CA-TOSG & $\tau=\TauStar$ & CA-TOSG & $\tau=\TauStar$ \\
\midrule
Scene-eq.\ $F_1$ & \TestRFFone & \TestTauFone & \CulverRFFone & \CulverTauFone \\
Payload (Msym)    & \TestRFPay  & \TestTauPay  & \CulverRFPay  & \CulverTauPay \\
$\Delta F_1$      & \multicolumn{2}{c}{\TestDeltaFone} & \multicolumn{2}{c}{\CulverDeltaFone} \\
$\Delta F_1$ LCB & \multicolumn{2}{c}{\TestLCB}       & \multicolumn{2}{c}{\CulverLCB} \\
Payload saving    & \multicolumn{2}{c}{\TestSaving\%}  & \multicolumn{2}{c}{\CulverSaving\%} \\
Scenes            & \multicolumn{2}{c}{\TestScenes}    & \multicolumn{2}{c}{\CulverScenes} \\
\bottomrule
\end{tabular}
\end{table}

\begin{figure*}[t]
\centering\includegraphics[width=\textwidth]{figures/fig2_primary}
\caption{Primary cell. Payload is drawn on a logarithmic axis and $\Delta F_1$ in its own panel:
the payload ratio is \TestRatio-fold while the $F_1$ shortfall on the bound is
\TestShortfall, so a shared scale would erase one half of the result. The right panel puts the preregistered margin and the bootstrap lower bound on one ruler.
On both splits the point estimate lies inside the margin and the bound does not.}
\label{fig:primary}
\end{figure*}

\subsection{Action mix}
The frozen selector requests E on \TestRhoEPct\% of Test rows and L on \TestRhoLPct\%, and never
requests F ($\rho_F=\TestRhoF$); the Culver-City mix is \CulverRhoEPct\% and \CulverRhoLPct\%. Section~\ref{sec:limits} explains why this is a property of the sampled
penalty grid rather than evidence that the feature action is useless.

\begin{figure}[t]
\centering\includegraphics[width=\columnwidth]{figures/fig3_recovery}
\caption{Delivery regimes on the development split. Fragment-aware partial recovery degrades
gracefully; all-or-nothing message delivery is governed by the survival probability $q$ rather than
by how the feature action performs.}
\label{fig:recovery}
\end{figure}

\subsection{Baselines}
Fixed-E, Fixed-L and Fixed-F bound the action space. The internal decision rules are the SNR
threshold and two- and three-scalar hand rules; the frozen $\tau=\TauStar$ is the primary comparator
because it is the only arm sharing CA-TOSG's complete payload definition. Where2comm~\cite{hu2022where2comm} is included as
an external feature-content baseline. It is evaluated under the same data splits, field of view,
ground truth and detection metrics. We report its native communication rate because the evaluated
implementation transmits floating-point selected features and does not execute the locked int8
quantisation path. Consequently, it is excluded from bit-level budget-matched claims. At a
communication threshold of $0.02$ it reaches AP@0.5 \WtwoCTestAPfifty{} on Test and
\WtwoCCulverAPfifty{} on Culver-City while retaining \WtwoCTestRate\% and \WtwoCCulverRate\% of
features. No adaptive-policy baseline met the reproducibility bar.

\begin{figure*}[t]
\centering\includegraphics[width=\textwidth]{figures/fig4_where2comm}
\caption{Where2comm on its native communication-rate axis. The horizontal axis is the fraction of
features retained, not a bit budget.}
\label{fig:w2c}
\end{figure*}

\section{Limitations}\label{sec:limits}
\textbf{Accuracy non-inferiority was not established.} On both splits the point estimate lies inside
the margin and the confidence bound does not. The margin was fixed in advance and the bound is what
it tests.

\textbf{Proportion.} The saving is overwhelming; the shortfall is small. Neither may be reported
without the other: stating only the saving packages a failed criterion as success, and stating only
the shortfall discards a \TestRatio-fold reduction.

Full-frame and collaborator-available payload savings differed by $0.000$ percentage points on both
held-out splits.

\textbf{Culver-City rests on \CulverScenes{} scenes.} The bootstrap resamples them, so its interval
is coarse; the result has the same shape as Test on far less scene diversity. Culver-City is a
\emph{simulated} domain shift.

\subsection{Why this remains a three-action method although $\rho_F=0$}
Four measurements belong together. First, F is genuinely the best action on $56.2\%$ of grid rows:
its realised utility exceeds both E and L there, so this is not an action that never helps. Second,
the break-even $\lambda$ falls in a gap in the sampled penalty grid---$0.01556$ conditional on the
rows where F is best, and $0.0048$ averaged over the whole grid, each value carrying its
conditioning set---while the grid samples $\lambda\in\{0,0.02,0.05,\dots\}$, so no candidate could
have been selected in the region where a mixed policy exists. Third, a finer $\lambda$ scan,
registered before it produced any number, shows a genuine three-action trade-off across that whole
interval, with $\rho_F$ falling smoothly from $0.534$ to $0.152$. Fourth, all of this is an
exploratory diagnostic on the development split: it enters no primary criterion, and the frozen
candidate was not replaced on the strength of it. Taken together, $\rho_F=0$ is a property of where
the penalty grid sampled, not evidence that the feature action is worthless.

\textbf{No three distinct budget points.} All three budgets select the same candidate and the budget
never binds. The contribution is therefore a granularity-control framework, not a demonstration that
a learned selector beats simple rules---at equal budget it does not.

\textbf{Other limits.} Selector-only latency is measured; total system latency is not. One
ego--collaborator pair. A flat block-fading channel.

\section{Conclusion}
Three transmission granularities were placed under one detector, one field of view and one ground
truth, and charged through one measured transport chain. Communication falls by a factor of
\TestRatio{} on Test and \CulverRatio{} on Culver-City relative to a budget-feasible threshold
policy, while accuracy non-inferiority is not established on either. The conditions are one dataset family, one
detector, a single ego--collaborator pair, and a flat block-fading channel.

\bibliographystyle{IEEEtran}
\bibliography{references}
\end{document}
